
\begin{proof}
Let $X =\{x \in \{0,1\}^{n/4}: |x| = \floor{n/8}-1\}$ and $Y =\{y \in \{0,1\}^{n/4}: |y| = \floor{n/8}+1\}$.  For every 
$x = (x^{(1)}, \ldots, x^{(3n/4)}) \in (X \cup Y)^{3n/4}$ we associate a bipartite graph $G_x = (L \sqcup R, w_x)$ 
where $L = \{1, \ldots, 3n/4\}, R = \{3n/4+1, \ldots, n\}$ and $w_x(\{i,j\}) = 1+ \varepsilon \cdot x^{(i)}(j-3n/4)$ for every $i \in L, j \in R$.  We set $\Gcal = \{G_x : x \in (X \cup Y)^{3n/4}\}$.

Define the function $g: X \cup Y \rightarrow \{0,1\}$ where $g(x) = 0$ iff $x \in X$.  We have $Q_{1/3}(g) \in \Omega(n)$ by \cref{thm:thresh}.  
Let $f : \{0,1\}^{3n/4} \rightarrow \{0,1\}$ be the AND function, for which $Q_{1/3}(f) \in \Omega(\sqrt{n})$.  By the composition theorem for quantum query complexity \cite{HLS07,Rei11a}, 
we have $Q_{1/3}(h) \in \Omega(n^{3/2})$ for the composed function $h = f \circ g^{3n/4}$.   

Let $x = (x^{(1)}, \ldots, x^{(3n/4)}) \in (X \cup Y)^{3n/4}$.  If $h(x) = 1$ then $x^{(i)} \in Y$ for all $i \in [3n/4]$ and the weight of the star cut $\Delta_{G_x}(\{i\}) = n/4 + \varepsilon \cdot(\floor{n/8}+1)$.  
As $\varepsilon \le 1$ this will be the weight of a minimum cut in $G_x$ by \cref{clm:complete_bipartite}.  On the other hand if $h(x) = 0$ then some $x^{(i)} \in X$ and 
$\Delta_{G_x}(\{i\}) = n/4 + \varepsilon \cdot(\floor{n/8}-1)$ and this will be the weight of a minimum cut of $G_x$.
Thus computing the weight of a minimum cut of $G_x$ lets us evaluate $h(x)$.  Further, given oracle access to $x$ we can 
simulate queries to $G_x$ in the adjacency array model.  Let $A_x$ be a $3n/4$-by-$n/4$ matrix whose $\ith$ row is the vector $1+ \varepsilon x^{(i)}$.  Then the vertical concatenation of $A$ with 
$A^T$ is a valid adjacency array for $G_x$.  To a degree query on vertex $i$ we simply answer $n/4$ if $1 \le 3n/4$ and $3n/4$ if $3n/4+1 \le i \le n$.  We can also answer 
a query to the name and weight of the $\jth$ neighbor of $i$ with one query to $x$.  This shows that the (1/3)-error quantum query complexity of computing the weight of 
a minimum cut on graphs in $\Gcal$ in the adjacency array model is at least $Q_{1/3}(f \circ g^{3n/4}) \in \Omega(n^{3/2})$.  

Finally, suppose a quantum query algorithm can compute a shore of a minimum cut in $G_x$ with $T$ queries.  We know that this shore must be of the form $\{v\}$ for a vertex $v \in L$.  
Thus with with $O(n)$ more queries the algorithm can classically compute the weight of a 
minimum cut by querying the weight of the neighbors of $v$.  Thus $T + O(n) \in \Omega(n^{3/2})$, which means $T \in \Omega(n^{3/2})$.  This completes the proof.
\end{proof}

\subsubsection{Large edge-weight ratio}
Let $n \in \N$ be a multiple of $4$ and $V$ be a vertex set with $|V| = n$.  Partition $V$ into four sets 
$V_1, V_2, V_3, V_4$.  

Now consider an integer $\tau$ such that $1 \leq \tau \leq 5n/8$ and $\tau n/10$ is an integer.  Fix a set $S$ of 
$\tau n/10$ ``edge disjoint'' quadruples $(u_1, u_2, u_3, u_4) \in V_1 \times V_2 \times V_3 \times V_4$.  By edge disjoint we mean no 
pair of consecutive elements $(u_i, u_{i+1})$ or $(u_4, u_1)$ appears in more than one quadruple. We fix an enumeration 
of $S$ and refer to the vertices in the $\ellth$ quadruple as $u_1^\ell, u_2^\ell, u_3^\ell, u_4^\ell$.  

For every $x \in \{0,1\}^{\tau n/10}$ we define an $n$-vertex weighted graph $G_x = (V,w_x)$ where 
$w_x(\{u,v\}) = \tau$ if $u \ne v \in V_i$ for some $i \in [4]$, and for $\ell \in [\tau n/10]$ we set
\begin{align*}
w_x(\{u^\ell_1,u^\ell_2\}) &= w_x(\{u^\ell_3,u^\ell_4\}) = x_\ell, \\
w_x(\{u^\ell_2,u^\ell_3\}) &= w_x(\{u^\ell_4,u^\ell_1\}) = 1 - x_\ell \enspace.
\end{align*}
Otherwise, $w_x(\{u,v\}) = 0$.  
In words, on each $V_i$ we have a complete graph with all edge weights $\tau$, and for each $\ell \in [\tau n/10]$ we either add unit weight edges
$\{u^\ell_1,u^\ell_2\}, \{u^\ell_3,u^\ell_4\}$ or $\{u^\ell_2,u^\ell_3\},\{u^\ell_4,u^\ell_1\}$ depending on $x_\ell$.  
The construction is depicted in \cref{fig:adj-lower-bound}.

There are a few important points to note about this definition.  First, the edge-weight ratio of $G_x$ is $\tau$.  Second, for any $X$ that nontrivially intersects some $V_i$ we have that 
$w(\Delta_{G_x}(X)) \ge \tau(n/4-1)$.  This means that such an $X$ cannot be the shore of a minimum cut of $G_x$.  Third, by construction 
the degree of every vertex of $G_x$ is independent of $x$.  This means that degree queries to $G_x$ can be trivially answered and give us 
no information about $x$.  

\begin{figure}[htb]
\centering
\includegraphics[width=.7\textwidth]{adj-lower-bound.pdf}
\caption{Figure of graph $G_x$.
If $x_\ell = 1$ then we add edges $\{u_1^\ell,u_2^\ell\}$ and $\{u_3^\ell,u_4^\ell\}$.
If $x_\ell = 0$ then we add edges $\{u_2^\ell,u_3^\ell\}$ and $\{u_4^\ell,u_1^\ell\}$.}
\label{fig:adj-lower-bound}
\end{figure}

\begin{lemma} \label{lem:adj-query}
We can simulate a single query to $G_x$ in the adjacency array model using a single query to $x$.
\end{lemma}
\begin{proof}
We first handle degree queries.  This can be answered with no queries to $x$ as the degree of a vertex is independent of $x$.

Now consider a query $(v,k) \in V \times [\deg(v)]$ to which we must answer the name $u$ of the $k$-th neighbor of 
$v$ and the edge weight $w_x(\{u,v\})$.  For clarity of exposition, we assume $v = u_1^t \in V_1$; the other cases are handled similarly.
\begin{itemize}
\item
If $k \leq n/4 - 1$ then return the $k$-th neighbor $u$ of $v$ inside $V_1$ and edge weight $w_x(\{u,v\}) = \tau$.
\item
If $k \ge n/4$ then let $j = k -n/4 +1$.  Letting $\ell$ denote the index of the $\jth$ quadruple of $S$ containing $v$ we query $x_\ell$.
\begin{itemize}
\item
If $x_\ell = 1$ then return neighbor $u^\ell_2$ and edge weight $w(\{v,u^\ell_2\}) = 1$.
\item
If $x_\ell = 0$ then return neighbor $u^\ell_4$ and edge weight $w(\{v,u^\ell_4\}) = 1$.
\end{itemize}
\end{itemize}
In total this takes a single query to $x$, which proves the lemma.
\end{proof}

\noindent
Now we can prove the following lemma.
\begin{lemma} \label{lem:adj-min-cut}
Fix integers $n$ and $\tau$ such that $1 \le \tau \le 5n/8$ and $\tau n/10 \in \N$.
Consider a string $x \in \{0,1\}^{\tau n/10}$ and the corresponding graph $G_x$.
If $|x| < \tau n/20$ then $G_x$ has a unique minimum cut with shores $(X,\overline X) = (V_1 \cup V_2, V_3 \cup V_4)$ and weight $w(\Delta_{G_x}(X)) = 2|x|$.
If $|x| > \tau n/20$ then $G_x$ has a unique minimum cut with shores $(X,\overline X) = (V_1 \cup V_4, V_2 \cup V_3)$ and weight $w(\Delta_{G_x}(X)) = 2(\tau n/10 - |x|)$.
\end{lemma}
\begin{proof}
First consider any cut shore that nontrivially intersects some $V_i$.
Since the subgraph $G_x[V_i]$ induced on $V_i$ is a complete graph with edge weights $\tau$, this implies that such a cut has weight at least $\tau (|V_i|-1) = \tau (n/4-1)$.
Now consider the small set of remaining cut shores that trivially intersect the $V_i$'s.
The weight of each one of these cuts can be easily expressed as a function of the Hamming weight $|x|$ of the input:
\begin{align*}
w(\Delta_{G_x}(V_i))
&= \tau n/10, \\
w(\Delta_{G_x}(V_1 \cup V_3)) &= w(\Delta_{G_x}(V_2 \cup V_4))
= 2 \tau n/10, \\
w(\Delta_{G_x}(V_1 \cup V_2)) &= w(\Delta_{G_x}(V_3 \cup V_4))
= 2 |x|, \\
w(\Delta_{G_x}(V_1 \cup V_4)) &= w(\Delta_{G_x}(V_2 \cup V_3))
= 2 (\tau n/10- |x|).
\end{align*}
It is clear that all minimum weight cuts will be among these cuts, and the lemma easily follows.
\end{proof}

Using this lemma we can prove the following theorem.

